\chapter{Simple Stellar Models and Compact Remnants}

\section{Core Stellar Structure Equations}

\subsection{Four Coupled ODEs}
\begin{theory}
For spherical stars:
\begin{align}
    \frac{dP}{dr} &= -\frac{Gm\rho}{r^2}, \\
    \frac{dm}{dr} &= 4\pi r^2\rho, \\
    \frac{dT}{dr} &= -\frac{3\kappa\rho}{16\pi acT^3}\frac{L(r)}{r^2}, \\
    \frac{dL}{dr} &= 4\pi r^2\rho\,\varepsilon.
\end{align}
Boundary conditions include $m(0)=0$, $L(0)=0$, and photospheric constraints near $r=R$.
\end{theory}

\subsection{Closure Relations}
\begin{foundation}
The structure equations require constitutive inputs:
\begin{itemize}
    \item Equation of state: $P=P(\rho,T)$
    \item Opacity: $\kappa=\kappa(\rho,T,\text{composition})$
    \item Nuclear source term: $\varepsilon=\varepsilon(\rho,T,\text{composition})$
\end{itemize}
\end{foundation}

\section{Equations of State and Opacity}

\subsection{Gas, Radiation, and Degeneracy Pressure}
\begin{theory}
\begin{align}
    P_{\text{gas}} &= \frac{\rho k_B T}{\mu m_H}, \\
    P_{\text{rad}} &= \frac{1}{3}aT^4, \\
    P_{\text{NR,deg}} &= K_{\text{NR}}\rho^{5/3}, \\
    P_{\text{UR,deg}} &= K_{\text{UR}}\rho^{4/3}.
\end{align}
Total pressure is the appropriate sum for the relevant regime.
\end{theory}

\subsection{Mean Molecular Weight}
\begin{theory}
For ionized composition fractions $(X,Y,Z)$,
\begin{equation}
    \frac{1}{\mu}\approx 2X+\frac{3}{4}Y+\frac{1}{2}Z.
\end{equation}
This parameter maps composition directly into the gas-pressure term.
\end{theory}

\subsection{Opacity and Nuclear Heating Scalings}
\begin{theory}
Useful approximations:
\begin{align}
    \kappa_{\text{es}} &\approx 0.02(1+X)\,\text{m}^2\text{kg}^{-1}, \\
    \kappa_{\text{bf+ff}} &\propto \rho T^{-3.5}, \\
    \varepsilon_{pp} &\propto \rho X^2 T^4.
\end{align}
\end{theory}

\begin{remark}
Opacity controls transport efficiency, while $\varepsilon(T)$ controls how centrally concentrated energy generation is.
Together they determine whether stellar zones are radiative or convective.
\end{remark}

\section{Analytic Scaling Models}

\subsection{Parameterized Density Profiles}
\begin{intuition}
Exact stellar models are numerical, but simple analytic profiles are valuable for quick estimates and physical intuition.
They expose which observables mostly constrain central density, pressure, and temperature.
\end{intuition}

\begin{theory}
A representative profile is
\begin{equation}
    \rho(r)=\rho_c\exp(-r^2/a^2)
\end{equation}
or polynomial alternatives.
With hydrostatic balance and ideal-gas EOS, one derives approximate $P(r)$ and $T(r)$ trends from the center to the surface.
\end{theory}

\section{Stellar Mass Limits}

\subsection{Minimum Mass for Hydrogen Burning}
\begin{theory}
Competition between thermal support and electron degeneracy yields a minimum mass:
\begin{equation}
    M_{\min}\approx 0.08\,M_\odot.
\end{equation}
Below this threshold, sustained H fusion does not ignite (brown dwarf regime).
\end{theory}

\subsection{Upper Mass and Radiation Support}
\begin{theory}
As mass increases, radiation pressure becomes significant, and instability / strong winds limit long-term growth.
Observed and theoretical upper masses are typically of order $100$--$150\,M_\odot$.
\end{theory}

\begin{importantnote}
The lower mass limit is set by degeneracy suppressing core heating;
the upper limit is set by radiation-driven instability and mass loss.
They arise from different physics and should not be conflated.
\end{importantnote}

\section{Compact Remnants}

\subsection{White Dwarfs}
\begin{theory}
White dwarfs are supported by electron degeneracy pressure, with characteristic behavior:
\begin{equation}
    R\propto M^{-1/3}\quad(\text{NR degenerate limit}),
\end{equation}
and maximum mass near
\begin{equation}
    M_{\text{Ch}}\approx 1.44\left(\frac{2}{\mu_e}\right)^2M_\odot.
\end{equation}
\end{theory}

\subsection{Cooling of White Dwarfs}
\begin{theory}
With ions as main heat reservoir and radiative surface losses,
\begin{equation}
    L\propto t^{-7/5}
\end{equation}
in simplified cooling models (modified by crystallization and envelope physics).
\end{theory}

\subsection{Neutron Stars (Overview)}
\begin{extra}
Neutron stars are relativistic compact objects requiring TOV hydrostatic balance:
\begin{equation}
    \frac{dP}{dr}=-\frac{G[\rho+P/c^2][m+4\pi r^3P/c^2]}{r^2\left[1-2Gm/(rc^2)\right]}.
\end{equation}
Observed masses around $2\,M_\odot$ strongly constrain the high-density EOS.
\end{extra}

\begin{example}{Regime check from mass scale}{ex:regime-check}
Compare two stars qualitatively:
\begin{itemize}
    \item $0.05\,M_\odot$: below H-burning threshold, degeneracy appears early.
    \item $30\,M_\odot$: strong radiation pressure and high CNO sensitivity.
\end{itemize}

\smallskip\textbf{Conclusion.}
Low-mass objects trend toward brown-dwarf behavior, while high-mass stars show luminous, unstable, short-lived evolution.
\end{example}
